\documentclass[10pt,a4paper]{amsart}
\usepackage[margin=19mm]{geometry}
\usepackage{amsmath,amssymb,mathtools}
\usepackage[scr=rsfs,cal=euler]{mathalfa}
\usepackage{xfrac}
\usepackage[unicode,hidelinks,bookmarks=false]{hyperref}
\newcommand{\ie}{i.e.\ }
\newcommand{\cf}{cf.\ }
\newcommand{\resp}{respectively\ }
\newcommand{\Subsection}{\subsection}
\providecommand{\nobreakdash}{\nobreak\hbox{-}}
\setlength{\emergencystretch}{2em}
\multlinegap0pt
\usepackage{bm}
\usepackage{bbm}
\usepackage{dsfont}
\usepackage{xspace}
\usepackage{cancel}
\usepackage{mathtools}
\usepackage{amsthm}
\makeatletter
\@ifundefined{theorem}{\newtheorem{theorem}{Theorem}}{}
\makeatother
\title[Explicit complex time integrators for stiff problems]{Explicit complex time integrators for stiff problems}
\author{Jithin D. George, Julian Koellermeier, Samuel Y. Jung,, Niall M. Mangan}
\date{Source: arXiv:2601.07730v1 (CC BY 4.0)}
\begin{document}
\maketitle

\noindent\emph{Source and licence.} Explicit complex time integrators for stiff problems. Version arXiv:2601.07730v1 (CC BY 4.0), distributed under the Creative Commons Attribution 4.0 International licence, as recorded in the arXiv licence metadata for this exact version and shown on the abstract page https://arxiv.org/abs/2601.07730v1. Full rights notice: licenses/arxiv-2601.07730-CC-BY-4.0.txt.\par\smallskip

\noindent\textit{Editorial scope.} Counted unit: the Theorem whose source label is \texttt{theorem1} in the article (source0.tex lines 313--315, source label \texttt{theorem1}), delivered together with the article's own complete original proof of it (source0.tex lines 317--365, one \texttt{proof} environment of 49 lines). No mathematics is written, repaired or re-derived: statement and proof are the article's own text line for line. Required context retained uncounted: none. Every definition, display and label of those lines is retained unchanged. Omitted from this package and not redistributed: 1--312, 316--316, 366--670. Nothing inside a retained excerpt is omitted or altered: every retained body is delivered exactly as the article's lines are written, so no omission receipt of any kind accompanies this package. Citations: the retained excerpts carry no citation command at all, so no bibliography entry of the article is transcribed and no reference is redistributed. Reference adaptation: none was needed and none was made; every reference inside the delivered unit resolves inside it, so no unresolved reference or citation is produced. Printed slips kept literal: no printed word, symbol, hypothesis or number of the counted statement or of its proof was altered. Layout and class: the article's own class is \texttt{siamonline190516} with packages including lipsum, amsfonts, graphicx, epstopdf, algorithmic, enumitem, amsopn, amsmath, mathtools, tabularx, epsfig, forest, cleveref; the wrapper is the lane's pinned \texttt{amsart} editorial wrapper at 10pt on A4, which restates the theorem-like environments the retained text needs and adds the article's own macro definitions that the retained text needs (none), with extra wrapper packages (bm, bbm, dsfont, xspace, cancel, mathtools). The delivered numbering is the unit's own. Assets: no retained span contains a figure environment, a graphics-inclusion command, a PDF-page inclusion, a table or a TikZ picture; no image file is copied into this package, the package declares no asset dependency and no image file is redistributed with it. Licence: version arXiv:2601.07730v1 of this article is distributed under the Creative Commons Attribution 4.0 International licence, as recorded in the arXiv licence metadata for this exact version and shown on the abstract page https://arxiv.org/abs/2601.07730v1. A full rights notice is delivered at licenses/arxiv-2601.07730-CC-BY-4.0.txt. Characters: every retained excerpt is pure ASCII, so the delivered engine, XeLaTeX with its default Latin Modern text font, reports no missing character.
\begin{theorem}\label{theorem1}
For a linear system $\dot{y} = A y$, where $y \in \mathbb{R}^n$ and $A \in \mathbb{R}^{n \times n}$ is a matrix whose eigenvalues $\lambda \in \mathbb{C} \setminus \mathbb{R}$ lie entirely on either the positive imaginary axis, i.e., $\lambda = +i|\lambda|$, or the negative imaginary axis, i.e., $\lambda = -i|\lambda|$, the optimal stability polynomial for a 2-stage first-order integrator has complex coefficients. 
\end{theorem}

\begin{proof}
Without loss of generality, consider a system $\dot{y} = A y$ with all the eigenvalues $\lambda$ of $A$ on the positive imaginary axis, i.e., $\lambda = +|\lambda|i$. 
The stability polynomial for a general 2-stage  first-order integrator is given by
\begin{align}
\label{eq:schro_stab}
    \Phi(z) = 1 + z + k z^2,
\end{align}
where the value of the coefficient $k$ needs to optimized so that the region enclosed by $|\Phi(z)| \leq 1$ is maximal on the positive imaginary axis, enabling the maximal stable time step.

For the stability polynomial with purely real coefficients, i.e., $k \in \mathbb{R}$, the stability condition $|\Phi(z)| \leq 1$    on $ z = iy \in [0, iy_{\text{max}}]$ becomes 
\begin{align}
\label{eq:phi2}
    %k^2y^4+(1-2k)y^2  = 
    \left(k^2y^2+(1-2k)\right) y^2 \leq 0,
\end{align}
where $y = +|\lambda| \Delta t$. For the integrator to be stable, we thus require $y^2 \in [0, \frac{2k-1}{k^2}$]. The maximal value of $y$ is 1, which is achieved for $k = 1$. Hence, for real-valued integrators, the optimal stability polynomial for this system is 

\begin{align}
    \Phi_r(z) = 1 +z + z^2
    \label{eq:stabschro_real}
\end{align} and the maximum stable step is $\Delta t = \frac{1}{|\lambda_{max}|}$.

For the stability polynomial with the complex coefficient $k \in \mathbb{C}$, $k = k_r + k_i i$, $k_r, k_i \in \mathbb{R}$, the stability condition translates to 
%\begin{align}
%\label{eq:phi1}
%    (k_r^2+k_i^2)y^4- 2 k_iy^3+ (1-2k_r)y^2 \leq 0
%\end{align}
%or 
\begin{align}
\label{eq:phi1b}
    (k_r^2+k_i^2)y^2- 2 k_iy+ (1-2k_r) =: f(y,k_r,k_i) \leq 0.
\end{align}
The stability needs to be fulfilled for vanishing $y=0$, requiring that the function $f(y,k_r,k_i)= (k_r^2+k_i^2)y^2- 2 k_iy+ (1-2k_r)$ needs to fulfill $f(0,k_r,k_i) \leq 0$, which fixes $k_r = 1/2$. The remaining upper root of $f(y,k_r,k_i)$ is given by $y = 2 k_i/(\frac{1}{4}+k_i^2)$, which attains a maximal value of $2$ for $k_i = \frac{1}{2}$. 

%Requiring the quadratic equation to be non-positive on an interval beginning at $y=0$ sets the constant term $(1-2k_r)$ to $0$ i.e  $k_r = 1/2$. The upper bound of the non-positive interval is $y = 2 k_i/(\frac{1}{4}+k_i^2)$, which achieves a maximal value of $2$ for $k_i = \frac{1}{2}$. 
 
Thus, the optimal coefficient $k =\frac{1}{2}+\frac{1}{2}i$ achieves a stability domain of $y \in [0, 2]$. Hence, the optimal stability polynomial for this system is 
\begin{align}
    \Phi_c(z) = 1 +z + \left(\frac{1}{2}+\frac{1}{2}i\right)z^2
\end{align}
 and the maximum stable step is $\frac{2}{|\lambda_{max}|}$ which is twice as large as its real equivalent.

Similarly, if all the eigenvalues lie on the negative imaginary axis like in the case of the linear Schrodinger equation, the optimal stability polynomial is 
\begin{align}
    \Phi_{c_-}(z) = 1 +z + \left(\frac{1}{2}-\frac{1}{2}i\right)z^2
    \label{eq:stabschro_complex}
\end{align}
and the maximum stable step is also $\frac{2}{|\lambda_{max}|}$.
\end{proof}

\end{document}
